\lipsum[5] % Dummy text

\begin{enumerate}[noitemsep] % [noitemsep] removes whitespace between the items for a compact look
\item First item in a list
\item Second item in a list
\item Third item in a list
\end{enumerate}

%------------------------------------------------

\subsection{Contact}

\lipsum[6] % Dummy text

\paragraph{Paragraph Description} \lipsum[7] % Dummy text

\paragraph{Different Paragraph Description} \lipsum[8] % Dummy text

%------------------------------------------------

\subsection{Cahier des charges fonctionnel}

\subsection{Cahier des charges Techniques}

\subsection{Outils et environnement de travail}

Pour réaliser ce projet, nous avons choisi des outils qui nous permettent de gérer toute la chaîne, du schéma électronique au code embarqué.

\subsubsection{Outils de collaboration} 
Dans un premier temps, nous avons mis en place un espace Google Drive \cite{Google_Drive} permettant à chaque membre du groupe de déposer son travail. Cependant, au vu de la nature technique du projet, nous nous sommes très vite dirigés vers GitHub \cite{GitHub}. Cet environnement s'est avéré beaucoup plus adapté pour la gestion de notre projet et le versioning, ce qui nous a permis de gagner un temps précieux. GitHub \cite{GitHub} nous a principalement permis d'avoir un espace de travail à plusieurs dans un environnement de programmation et de schématisation, que ce soit pour KiCad ou VS Code, ce qui n'était pas permis par Google Drive \cite{Google_Drive}. L'arborescence de notre espace de travail ainsi que la gestion de nos différents fichiers (rapport, codes, schémas) sont illustrées sur la Figure \ref{fig:GitHub}. 
\begin{itemize} 
    \item \textit{Historique :} Nous enregistrons chaque modification. Pour le debug, nous pouvons facilement identifier quelle est la dernière version fonctionelle et qui a effectué quel changement et quand. 
    \item \textit{Travail en groupe :} Nous pouvons avancer sur l'électronique, le code ou la mécanique en parallèle sans nous gêner. Mais aussi travailler sur le même fichier depuis deux ordinateurs. 
    \item \textit{Sauvegarde :} Tout est stocké et versionné, ce qui nous permet de ne pas perdre nos fichier accidentellement, et de régulièrement avoir des points de sauvegarde dans l'avancement du projet. 
\end{itemize} 

\begin{figure}[!ht]
\centering
\includegraphics[width=0.8\columnwidth]{github_gestion_projet.png}
\caption[Interface du dépôt GitHub utilisé pour le travail collaboratif]{Interface du dépôt GitHub utilisé pour le travail collaboratif}
\label{fig:GitHub}
\end{figure}



\subsubsection{Développement et code}
\begin{itemize}
    \item \textbf{VS Code :} C'est l'éditeur que nous utilisons pour tout. Comme nous le maîtrisons tous, nous avons pu garder nos habitudes de travail tout en centralisant nos tâches. Il nous sert pour le code, les commandes Git et le développement matériel.
    \item \textbf{PlatformIO :} Intégré à VS Code, c'est l'outil que nous utilisons pour programmer nos microcontrôleurs.
    \begin{itemize}
        \item \textit{Gestion multi-projets :} Nous passons d'une configuration à l'autre (Hub ou ancres) via les fichiers \texttt{platformio.ini}, ce qui nous évite des erreurs.
        \item \textit{Bibliothèques :} Elles sont téléchargées et mises à jour automatiquement.
        \item \textit{Git :} Nous voyons les infos de version directement dans l'interface, ce qui nous permet de suivre nos modifications en temps réel.
    \end{itemize}
\end{itemize}

\subsubsection{Conception électronique et mécanique}
\begin{itemize}
    \item \textbf{KiCad :} Nous l'utilisons pour toute la partie électronique. Nous avons créé nos propres composants, routé nos cartes et préparé les fichiers de fabrication. Cela nous a permis d'intégrer précisément des modules comme le BMP581 ou le module UWB.
    \item \textbf{Autodesk Fusion 360 :} Pour la partie 3D, nous importons le modèle STEP depuis KiCad pour concevoir le boîtier autour de la carte. Nous calons ainsi les fixations et les ouvertures au millimètre.
    \item \textbf{Bambu Studio :} Pour l'impression 3D, nous l'utilisons pour configurer nos pièces, vérifier la structure avant impression et lancer la fabrication sur les machines de l'ESIEE.
\end{itemize}

